\section{Problem Formulation}
\label{sec:formulation}

The introduction established the empirical puzzle: parameter trajectories under signal transformations are not arbitrary.
We now formalize the central objects of our analysis.
Our starting point is the observation that a transformed signal $T_g[f]$ has its own optimal parameter vector $\theta^*(T_g[f])$, and the set of all such pairs traces a curve in weight space.
Understanding the shape of this curve---its curvature, its dimension, its global topology---is the goal of everything that follows.

\subsection{Optimal Parameter Orbit under Group Transformations}

Let $f: \Omega \to \R^d$ be a target signal (e.g., an image) and let $G$ be a transformation group (e.g., $\mathbb{R}^+$ for scale, $\SO(2)$ for rotation) acting on the signal space via $T_g \cdot f$.
We define:

\begin{definition}[Optimal Parameter Orbit]
The optimal parameter orbit of $f$ under $G$ is the set
\begin{equation}
    \mc{O}(f) = \{\theta^*(T_g \cdot f) \mid g \in G\} \subset \Theta,
\end{equation}
where $\theta^*(f) = \arg\min_{\theta} \mc{L}(f, f_\theta)$ is the global minimizer of a reconstruction loss.
\end{definition}

Since INR is a highly nonlinear composite mapping, a global linear equivariance assumption $\theta^*(T_g[f]) \approx \rho(g)\theta^*(f)$ lacks prior justification.
We instead study the \textbf{local tangent space geometry} of this orbit.

\subsection{Local Induced Tangent Map}

Let $\mathfrak{g}$ be the Lie algebra of $G$.
For $X \in \mathfrak{g}$, define the orbit tangent vector:

\begin{definition}[Local Induced Tangent Map]
\begin{equation}
    J_X(f) := \left.\frac{d}{dt} \theta^*(T_{\exp(tX)}[f])\right|_{t=0} \in T_{\theta^*(f)}\Theta.
    \label{eq:tangent_map}
\end{equation}
\end{definition}

When the symmetry error
\begin{equation}
    \eps_{\mathrm{sym}}(X) := \frac{\|\theta^*(T_{\exp(X)}[f]) - \mathrm{Exp}_{\theta^*(f)}(J_X(f))\|}{\|\theta^*(f)\|}
    \label{eq:symmetry_error}
\end{equation}
is sufficiently small, $J_X(f)$ approximately satisfies the Lie algebra homomorphism condition:

\begin{equation}
    [J_X, J_Y] \approx J_{[X,Y]}.
    \label{eq:lie_homomorphism}
\end{equation}

In this regime, $\{J_X(f)\}_{X \in \mathfrak{g}}$ constitutes an \textbf{approximately linear representation} of $\mathfrak{g}$ in the tangent space $T_{\theta^*(f)}\Theta$.

\subsection{Spectral Structure and Group Type}

The spectral decomposition of $J_X$ directly determines the orbit geometry:

\begin{prop}[Group Type Dictates Orbit Geometry]
Let $J_X$ be the tangent map for a one-parameter subgroup $\exp(tX)$.

\begin{itemize}
    \item \textbf{Scale group $\mathbb{R}^+$ (non-compact)}: The Lie algebra $\mathfrak{g}_s \cong \mathbb{R}$ has real eigenvalues.
    Predicted: Parameter trajectories extend monotonically along exponential manifolds (open orbit).

    \item \textbf{Rotation group $\SO(2)$ (compact)}: The Lie algebra $\mathfrak{so}(2) \cong \mathbb{R}$ has purely imaginary eigenvalues $\pm ik$.
    Predicted: Trajectories are elliptical in the local 2D rotation subspace and close globally ($360^\circ$ returns to origin).
\end{itemize}
\end{prop}

This derivation strictly avoids global linear assumptions and is fully compatible with the NTK local linearization framework.

\input{figures/fig_tangent_space_decomposition}

\subsection{Research Questions}

We organize our investigation around three progressive questions:

\begin{itemize}
    \item[\textbf{A.}] \textbf{Existence}: Do parameter orbits $\mc{O}(f)$ exhibit detectable low-dimensional structure under signal group transformations?

    \item[\textbf{B.}] \textbf{Structural bounds}: Does the tangent map $J_X$ approximately satisfy Lie algebra homomorphism, and can we quantify perturbation stability?

    \item[\textbf{C.}] \textbf{Fast adaptation application}: Can we leverage the true orbital tangents for geometric-aware initialization in conditional INR?
\end{itemize}
